\begin{tikzpicture}[x=1cm,y=1cm,state/.style={draw,rounded corners,fill=Pale,note,minimum width=2.1cm,minimum height=.8cm}]
\path[use as bounding box](-.3,-.7) rectangle (15.5,4.9);
\node[note,anchor=west,font=\bfseries\sffamily]at(0,4.5){Joining branches / four states, four labeled edges};
\node[state](s)at(1.8,2.2){$A+B+C$};
\node[state](u)at(7.4,3.4){$AB+C$};
\node[state](v)at(7.4,1){$AC+B$};
\node[state,fill=Green!10](t)at(13.2,2.2){$ABC$};
\draw[flow](s)--node[above,sloped,smallnote]{bind AB}(u);
\draw[flow](s)--node[below,sloped,smallnote]{bind AC}(v);
\draw[flow](u)--node[above,sloped,smallnote]{bind C}(t);
\draw[flow](v)--node[below,sloped,smallnote]{bind B}(t);
\node[smallnote,text width=6cm] at (3.8,-.1){One canonical final vertex; no loss of alternative reaction labels.};
\node[smallnote,text width=6cm] at (11.4,-.1){First-discovery predecessor selects one shortest witness; the edge set preserves both routes.};
\end{tikzpicture}
